\section{Distance-\texorpdfstring{$1$}{1} domination}\label{sec:distone}

In this section we give simpler data structures for the distance-$1$ case, that is, implementing queries $\near{S,1}$ and $\far{S,1}$ under the assumption that the maintained graph is $d$-degenerate. These data structures are new; they are not merely specializations of the bounded-expansion constructions from \cref{sec:implementation}. 

Throughout this section we assume that the maintained dynamic graph $G$ has a fixed vertex set $V=\{1,\ldots,n\}$ and is $d$-degenerate at all times, for a fixed parameter $d\in \N$. We also assume, without loss of generality, that $n\geq 2$ and $d\geq 1$.


% and, within the constructions below, assume $n\ge 2$ and $d\ge 1$. The exceptional cases needed for \cref{thm:main-domset-deg} are handled explicitly in the proof of that theorem at the end of this section.

%In this section we restate Theorem~\ref{thm:NearVertex-new} and Theorem~\ref{thm:FarVertex-new} after specializing them to distance $r=1$ and to dynamic $d$-degenerate graphs on a fixed vertex set $V=\{1,\ldots,n\}$, where $n \ge 2$ and $d\ge 1$.
%\bbinline{Obecne pierwsze zdanie jest mylące: nie jest to przeformułowanie konstrukcji dla bounded expansion, lecz nowa i prostsza konstrukcja dla odległości $1$ w grafach $d$-zdegenerowanych. Proponuję zastąpić pierwszy akapit tej sekcji następującym tekstem po angielsku:

%``In this section we give simpler data structures for the distance-$1$ case under the assumption that the maintained graph is $d$-degenerate. These data structures are new; they are not merely specializations of the bounded-expansion constructions from \cref{sec:implementation}. We work on a fixed vertex set $V=\{1,\ldots,n\}$ and, within the constructions below, assume $n\ge 2$ and $d\ge 1$. The exceptional cases needed for \cref{thm:main-domset-deg} are handled explicitly in the proof of that theorem at the end of this section.''}

Using the Brodal--Fagerberg data structure (\cref{thm:bf}), we may assume that we maintain an orientation $\vec{G}$ of $G$ with maximum outdegree $\Delta \coloneqq 4d$. 
We write $N(v)$ for the open neighborhood of $v$ in $G$ (consisting of all the neighbors of $v$ in $G$), $N[v]\coloneqq \{v\}\cup N(v)$ for the closed neighborhood, and $N[X]\coloneqq \bigcup_{v\in X}N[v]$ for any $X\subseteq V$. We also write $N^+(v)\coloneqq \{w\in V : (v,w)\in E(\vec G)\}$ for the open outneighborhood of $v$ in $\vec G$.

\begin{comment}
\subsection{Arboricity and dynamic orientation}\label{sec:distone-arboricity}

A $d$-degenerate graph has arboricity at most $d$.

We use the following result of Brodal and Fagerberg~\cite[Theorem~3]{BrodalF99}.

\begin{theorem}[Brodal--Fagerberg]\label{thm:bf-bartek}
There exists a deterministic data structure that maintains a graph of arboricity at most $\alpha$ under edge insertions and deletions.
The data structure stores an orientation $\vec{G}$ of $G$ with maximum out-degree $4\alpha$, i.e., for every vertex $u$ we have $|N^+(u)| \le 4\alpha$.
An edge insertion has amortized cost $\Oh(1)$, an edge deletion has amortized cost $\Oh(\alpha + \log n)$, and the amortized number of edges that change orientation in a single operation is $\Oh(\alpha + \log n)$.
The data structure uses $\Oh(n + m)$ memory.
\end{theorem}

\begin{lemma}[Reporting reoriented edges]\label{lem:bf-report}
There exists an implementation of the data structure from Theorem~\ref{thm:bf} such that each update additionally outputs the list of edges whose orientation changed during this update.
If $r$ is the number of edges reoriented during this update, then producing this list takes $\Oh(r)$ additional time.
\end{lemma}
\begin{proof}
We augment the update procedure as follows.
Whenever the implementation changes the orientation of an edge, it appends this edge to a temporary list associated with the current update.
Thus each reoriented edge contributes $\Oh(1)$ additional work and is recorded exactly once, so producing the whole list takes $\Oh(r)$ time.
\end{proof}
\end{comment}

\subsection{Toolkit}\label{sec:distone-infra}

%Throughout this section we assume $n\ge 2$ and $d\ge 1$.
%\bbinline{W tej sekcji zakładamy $n\ge 2$ i $d\ge 1$, natomiast \cref{thm:main-domset-deg} dopuszcza także przypadki $k=0$, $n\le 1$ i $d=0$. Proponuję wstawić poniższy akapit jako pierwszy akapit dowodu \cref{thm:main-domset-deg}:

%``We first dispose of the degenerate parameter cases. If $k=0$, return $\emptyset$ if and only if $V(G)=\emptyset$, and return $\bot$ otherwise. If $n\le 1$ or $d=0$, then $G$ is edgeless at all times; in this case return $V(G)$ if $|V(G)|\le k$, and return $\bot$ otherwise. Hence, for the remainder of the proof, we may assume that $n\ge 2$, $d\ge 1$, and $k\ge 1$.''}
%By the arboricity bound above, Theorem~\ref{thm:bf}, and %Lemma~\ref{lem:bf-report}, the current $d$-degenerate graph admits a maintained orientation $\vec G$ of maximum out-degree $\Delta := 4d$.
%We write .

We start with describing a helper data structure that we call \emph{toolkit}, which will be later used in the data structures for $\near{S,1}$ and $\far{S,1}$ queries. We stress that the toolkit is fully deterministic; randomization will enter the scene only later, through an application of \cref{thm:dvorak-tuma-examples}.

%This subsection describes one deterministic toolkit for distance $1$ in dynamic $d$-degenerate graphs.
%Later data structures instantiate this toolkit independently; the description is common only at the level of exposition.

Let us first explain the intuition.
The toolkit architecture replaces searching through neighborhoods in $G$ (which may be large) with certain bucket arithmetic plus inspection of a few outneighborhoods in~$\vec G$; these are small due to the bound on the maximum outdegree of $\vec G$.  %guarantees that every set $N^+(u)$ has size at most $\Delta$.
Buckets reduce the query ``given $Z$, find $u$ such that $Z\subseteq N^+(u)$'' to a single dictionary lookup: instead of scanning through all the vertices of~$V$, we consult the single list $B(Z)\coloneqq \{u\in V \mid Z\subseteq N^+(u)\}$, maintained in the dictionary.

To facilitate later uses for weighted counting,
we shall assume that the maintained graph is equipped with a weight function $\wei\colon V\to \Z$ that is fixed upon initialization and does not change over time. In the toolkit data structure, we then explicitly maintain the following objects at all times.
\begin{itemize}
%\item The set $V^+\coloneqq \{u\in V \mid \deg_G(u)>0\}$, represented by a Boolean array indexed by $V$, an iterable list of its current elements, and its current cardinality.
%\item The set of isolated vertices $V\setminus V^+$ as a dictionary of disjoint maximal identifier intervals; this supports finding the smallest isolated vertex, or the next one after a given identifier, in $\Oh(\log n)$~time.
\item The neighbor dictionaries $N(u)$ and $N^+(u)$ for every vertex $u\in V$.
\item The bucket family $B(Z)\coloneqq \{u\in V \mid Z\subseteq N^+(u)\}$ for all sets $Z\subseteq V$ with $B(Z)$ being nonempty (note that this implies $|Z|\leq \Delta$). Each bucket $B(Z)$ is stored as a linked list, while the whole bucket family is stored as a dictionary keyed by the sorted representation of $Z$. Observe that every vertex $u\in V$ is on $2^{|N^+(u)|}\leq 2^\Delta$, hence the total sum of the lengths of the lists stored for the buckets $B(Z)$ is bounded by $2^\Delta\cdot n$. In particular, the dictionary stores at most this many entries. Additionally, with every bucket $B(Z)$ we store its total weight $\wei(B(Z))\coloneqq \sum_{v\in V(Z)}\wei(v)$.
\item A registration list for every vertex $u \in V$, tracking its $2^{|N^+(u)|} \le 2^\Delta$ bucket memberships (pointers to list elements on the bucket lists).
\end{itemize}
For every vertex $u$, the dictionaries $N(u)$ and $N^+(u)$ are keyed by vertex identifiers.
Hence for every $v\in V$, the membership tests $v\in N(u)$ and $v\in N^+(u)$ can be performed in $\Oh(\log n)$ time.

%As mentioned, only nonempty buckets are materialized in the bucket dictionary.
Every bucket key $Z$ is stored explicitly as the sorted list of its elements, so it occupies $\Oh(|Z|)\le \Oh(\Delta)=\Oh(d)$ space.
Every maintained bucket stores a pointer to its first element, its current cardinality, and its current weight.
Consequently, one bucket lookup, insertion, or deletion by key costs $\Oh(\log (2^\Delta n))=\Oh(\Delta+\log n)$, while after the bucket is found, accessing its first element or its cardinality or its weight costs $\Oh(1)$ time. Note that the bucket $B(\emptyset)$ consists of the whole vertex set $V$ at all times.
% checking whether a vertex belongs to $V^+$ costs $\Oh(1)$, and reading $|B(Z)|$ or $|V^+|$ costs $\Oh(1)$.
% time by setting $V^+=\emptyset$, the isolated-vertex dictionary to the single interval $[1,n]$, and all neighbor and bucket dictionaries to be~empty.

%For later use in weighted structures, every update of one toolkit instance additionally outputs a transient log consisting of:
%\begin{itemize}
%\item the list $L_{\mathrm{in}}^+$ of vertices that entered $V^+$;
%\item the list $L_{\mathrm{out}}^+$ of vertices that left $V^+$;
%\item the list $\Lambda^+$ of added bucket registrations $(u,Z)$;
%\item the list $\Lambda^-$ of deleted bucket registrations $(u,Z)$.
%\end{itemize}
%The proof below shows that these items are already traversed by the update routine, and in particular every logged registration event is emitted while the corresponding bucket object is already identified, so emitting the log does not change the asymptotic complexity.

We now prove a lemma providing guarantees on the complexity of maintaining a toolkit.

\begin{lemma}[Toolkit]\label{lem:infra-update}
One instance of the toolkit described in this subsection can be initialized on the edgeless graph on $V$ provided with a weight function $\wei\colon V\to \Z$ in $\Oh(n)$ time.
The toolkit uses $\Oh(2^{4d}\cdot dn)$ space at all times.
Upon edge insertion or deletion, the toolkit can be updated in amortized time $\Oh(2^{4d}\cdot d^2\cdot\log^2 n)$. % and outputs the complete log $(L_{\mathrm{in}}^+,L_{\mathrm{out}}^+,\Lambda^+,\Lambda^-)$.
%The total number of logged items is $\Oh(2^{4d}\cdot d\cdot \log n)$ in the amortized sense.
\end{lemma}
\begin{proof}
Initialization takes $\Oh(n)$ time: we set up all the relevant dictionaries to be empty, and then register all the vertices in the bucket $B(\emptyset)$, which becomes the only element in the dictionary of buckets.

%Initialization is exactly the constant-time setup described above, so only updates require analysis.

Consider one edge insertion or deletion.
First, we update the two affected neighbor dictionaries $N(\cdot)$; this takes $\Oh(\log n)$ time.
%, the Boolean flags for membership in $V^+$, the iterable list of $V^+$, and the interval dictionary for $V\setminus V^+$.
%Each endpoint can change its isolated/non-isolated status at most once during this step, so this first stage costs $\Oh(\log n)$ time.
%Whenever an endpoint enters $V^+$ we append it to $L_{\mathrm{in}}^+$, and whenever an endpoint leaves $V^+$ we append it to $L_{\mathrm{out}}^+$.
%Thus the total size of these two lists is $\Oh(1)$.
Second, we update the Brodal--Fagerberg orientation $\vec G$.
By \cref{thm:bf}, this takes amortized time $\Oh(d+\log n)$, the amortized number of reoriented edges is $\Oh(\log n)$, and we obtain the whole list of those edges, hence also the set of vertices whose out-neighborhood changed.
Each reoriented edge contributes at most its two endpoints.
In addition, the updated edge contributes only its two endpoints. %, because only these endpoints can change their membership in $V^+$.
Therefore, the amortized number of vertices $u$ for which the set $N^+(u)$ may change is~$\Oh(\log n)$.

Fix one such vertex $u$.
Let $A_{\mathrm{old}}\coloneqq N^+_{\mathrm{old}}(u)$ and $A_{\mathrm{new}}\coloneqq N^+_{\mathrm{new}}(u)$, where we follow the convention that $\mathrm{old}$ and $\mathrm{new}$ signifies whether an object in $G$ refers to the state before or after the update.
The old registrations of $u$ form the family $\mathcal{R}_{\mathrm{old}}(u)\coloneqq \{Z\mid Z\subseteq A_{\mathrm{old}}\}$, while the new registrations form the family  $\mathcal{R}_{\mathrm{new}}(u)\coloneqq \{Z\mid Z\subseteq A_{\mathrm{new}}\}$.
%And if $u\notin V^+_{\mathrm{old}}$, then $\mathcal{R}_{\mathrm{old}}(u)=\emptyset$, and if $u\notin V^+_{\mathrm{new}}$, then $\mathcal{R}_{\mathrm{new}}(u)=\emptyset$.
%In particular, if an endpoint becomes isolated after deleting its last incident edge, then it is inserted into no bucket after the update, not even into $B(\emptyset)$.
%Likewise, if a previously isolated endpoint becomes non-isolated, then it had no old registrations and acquires exactly the registrations indexed by subsets of $A_{\mathrm{new}}$.
%Since $|A_{\mathrm{old}}|,|A_{\mathrm{new}}|\le \Delta$, both families have size at most $2^\Delta$.
%We enumerate all sets in $\mathcal{R}_{\mathrm{old}}(u)$ and delete the corresponding registrations, then enumerate all sets in $\mathcal{R}_{\mathrm{new}}(u)$ and insert the new registrations.
Producing one subset key and performing the corresponding bucket lookup, insertion or deletion, and update to the bucket's cardinality and weight, costs $\Oh(\Delta+\log n)$, because the key has size at most $\Delta$. Note here that if the bucket $B(Z)$ did not exist before ($B_{\textrm{old}}(Z)=\emptyset$), it must be added to the dictionary of buckets, and if its cardinality drops to $0$, then it must be removed from the dictionary of buckets; all of this can be done within the same complexity bounds.
So the whole work caused by the vertex $u$ is bounded by \[\Oh((|\mathcal{R}_{\mathrm{old}}(u)|+|\mathcal{R}_{\mathrm{new}}(u)|)\cdot (\Delta+\log n))=\Oh(2^\Delta\cdot (\Delta+\log n)).\]
%Whenever one old registration $(u,Z)$ is deleted, we append $(u,Z)$ to $\Lambda^-$, and whenever one new registration $(u,Z)$ is inserted, we append $(u,Z)$ to $\Lambda^+$.
%Recording these events costs $\Oh(1)$ per registration,  and gives rise to $\Oh(2^\Delta\cdot \Delta)$ records in the log.

Multiplying by the amortized number $\Oh(\log n)$ of changed vertices per graph update yields amortized update time $\Oh(2^\Delta\cdot \Delta\cdot (\Delta+\log n)\cdot \log n)\leq \Oh(2^{4d}\cdot d^2\cdot \log^2 n)$, because $\Delta=4d$.
Note here that the factor $2^{4d}$ comes from enumerating all subsets of an out-neighborhood of size at most $\Delta=4d$.
%The same calculation shows that the combined size of $\Lambda^+$ and $\Lambda^-$ is $\Oh(2^\Delta \cdot \Delta\cdot \log n)=\Oh(2^{4d}\cdot d\cdot \log n)$ in the amortized sense.
%Together with the two lists of vertices entering or leaving $V^+$, this proves the bound on the update log.

For the persistent space bound, the orientation $\vec G$ and the (out)neighbor dictionaries use $\Oh(nd)$ words because every $d$-degenerate graph on $V$ has at most $d(n-1)$ edges.
For every $u\in V$, its registration list contains at most $2^{|N^+(u)|}\le 2^\Delta$ entries.
Each such entry is associated with a bucket key $Z\subseteq N^+(u)$, and if this key is stored as a sorted list of identifiers, then it uses $\Oh(|Z|)\le \Oh(\Delta)$ words.
Hence all registrations and bucket incidences together use at most $\Oh\left(\sum_{u\in V} 2^{|N^+(u)|}\cdot \Delta\right)=\Oh(2^\Delta\cdot \Delta\cdot n)$ words, which also gives an upper bound on the total size of the lists stored for the buckets. 
%The number of materialized buckets is at most the number of nonempty registrations plus one, so their headers, including the stored keys, satisfy the same bound.
Thus, as $\Delta=4d$, the whole toolkit uses $\Oh(2^{4d}\cdot dn)$ words of space.
\end{proof}

In further preparation for the queries $\near{S,1}$ and $\far{S,1}$, we introduce the following defintions. Let $S\subseteq V(G)$. We define:
\begin{align*}
	U(S)&\coloneqq \{u\in V\setminus S\mid S\subseteq N(u)\};\\
	U_2(S) & \coloneqq \{u\in V\setminus S\mid S\subseteq N^+(u)\};\\
	U_1(S) & \coloneqq U(S)\setminus U_2(S)=\{u\in U(S)\mid u\in \textstyle\bigcup_{s\in S} N^+(s)\}.
\end{align*}
Clearly, $(U_1(S),U_2(S))$ is a partition of $U(S)$.
Note also that $|U_1(S)|\leq \Delta|S|$, because $|N^+(s)|\leq \Delta$ for each $s\in S$.
Finally, observe the following.

%\begin{claim}[Decomposition of $U(S)$]\label{lem:distone-U1-bound}
%Assume $n\ge 2$.
%Let $S\subseteq V$ be nonempty and define $U(S):=\{u\in V\setminus S : S\subseteq N(u)\}$.
%Let $U_2(S):=\{u\in U(S) : S\subseteq N^+(u)\}$ and $U_1(S):=U(S)\setminus U_2(S)$.
%Then $U(S)=U_1(S)\cup U_2(S)$ is a disjoint union %and $U_1(S)\subseteq \bigcup_{s\in S}N^+(s)$.
%In particular, $|U_1(S)|\le \Delta\cdot |S|$.
%\end{claim}
%\begin{claimproof}[Proof.]
%The decomposition $U(S)=U_1(S)\cup U_2(S)$ is immediate from the definitions.
%Let $u\in U_1(S)$.
%Then $S\subseteq N(u)$ but $S\not\subseteq N^+(u)$, so there exists $s\in S$ with $s\notin N^+(u)$.
%Since $\{u,s\}\in E$ and the edge is not oriented out of $u$, it must be oriented as $(s,u)$.
%Hence $u\in N^+(s)$.
%\end{claimproof}

\begin{lemma}[Bucket characterization of $U_2(S)$]\label{lem:distone-U2-bucket}
For any $S\subseteq V$, we have $U_2(S)=B(S)$.
\end{lemma}
\begin{proof}
	Follows from the definitions and observation that $S\subseteq N^+(u)$ entails $u\notin S$, due the fact that $G$ is a simple graph.
\end{proof}
%\begin{proof}
%If $|S|>\Delta$, then $U_2(S)=\emptyset$ because every outneighborhood has size at most $\Delta$.
%Suppose then that $|S|\le \Delta$.
%Let $u\in U_2(S)$.
%Then $u\in U(S)$ and $S\subseteq N^+(u)$, so $u\in B(S)$.
%Conversely, let $u\in B(S)$.
%Then $u\in V^+$ and $S\subseteq N^+(u)\subseteq N(u)$.
%Since $S\subseteq N^+(u)$ and a vertex cannot be an out-neighbor of itself in a simple graph, we have $u\notin S$.
%Therefore $u\in U(S)$, and together with $S\subseteq N^+(u)$ this yields $u\in U_2(S)$.
%\end{proof}

\cref{lem:distone-U2-bucket} provides us with the following combinatorial characterization of vertices that are candidates for the output in the query $\near{S,1}$.

\begin{lemma}\label{cl:distone-near-candidates}
Let $S\subseteq V$ be nonempty, and suppose every vertex of $S$ is non-isolated in $G$.
If a vertex $v$ satisfies $S\subseteq N[v]$, then one of the following holds:
\begin{itemize}
\item $v\in S$; or
\item $v\in \bigcup_{s\in S}N^+(s)$; or
\item $|S|\le \Delta$ and $v\in B(S)$.
\end{itemize}
In particular, apart from the single bucket $B(S)$, there list above includes at most $|S|+\Delta\cdot |S|=\Oh(d|S|)$ candidates.
\end{lemma}
\begin{proof}
If $v\in S$, we are in the first case.
Assume then that $v\notin S$, so $v\in U(S)$.
If $v\in U_1(S)$, then $v\in \bigcup_{s\in S}N^+(s)$.
And if $v\in U_2(S)$, then $|S|\leq \Delta$ due to $S\subseteq N^+(v)$, and \cref{lem:distone-U2-bucket} yields $v\in B(S)$.
\end{proof}

\subsection{Data structure for \texorpdfstring{$\near{S,1}$}{near(S,1)}}\label{sec:distone-near}

%\paragraph{The data structure \texorpdfstring{$\nearVertexDS{G}{k,1}$}{NearVertex[1,n]}.}

In this subsection we provide the desired implementation of the $\near{S,1}$ queries. Note that this time, the data structure is entirely deterministic and can handle any given set $S$, not necessarily of size stipulated by a fixed constant. Also, since we are going to reuse this data structure in the implementation of $\far{S,1}$, we equip it with the capability of weighted counting.

\begin{theorem}\label{thm:NearVertex-new-deg}
	Fix $d\in \N$.
	Let $G$ be a dynamic graph on vertex set $V=\{1,\ldots,n\}$ that is $d$-degenerate at all times, and let $\wei\colon V\to \Z$ be a weight function fixed upon initialization. Then there exists adata structure $\nearVertexDS{G}{}$ that provides access to the following query:
	\begin{itemize}
		\item $\near{S,1}$: For a given set of vertices $S\subseteq V(G)$ with $|S|\leq k$, return a vertex $v$ such that $S\subseteq N[v]$, or $\bot$ if no such vertex exists, as well as the value $\sum_{v\colon S\subseteq N[v]} \wei(v)$.
	\end{itemize}
	The amortized time of an update is $\Oh(2^{4d}\cdot d^2\cdot \log^2 n)$. The queries take worst-case time $\Oh(d|S|^2\cdot\log n)$. The data structure can be initialized for an edgeless $G$ in $\Oh(n)$ time and uses $\Oh(2^{4d}\cdot dn)$ space at all times.
\end{theorem}
\begin{comment}
The structure $\nearVertexDS{G}{k,1}$ consists of one instance of the toolkit of \ref{sec:distone-infra}; it stores no additional persistent fields.

\bbinline{Obecna notacja $\nearVertexDS{G}{k,1}$ i $\farVertexDS{G}{k,1}$ odwraca kolejność parametrów względem \cref{lem:summary-dom}; ponadto parametr $k$ nie ogranicza tu dopuszczalnego rozmiaru wejściowego zbioru. Proponuję w całej tej sekcji wykonać dosłowne zamiany
\[
 \nearVertexDS{G}{k,1}\quad\text{przez}\quad\nearVertexDS{G}{1},
 \qquad
 \farVertexDS{G}{k,1}\quad\text{przez}\quad\farVertexDS{G}{1}.
\]
Następnie proponuję wstawić po pierwszym akapicie wprowadzającym strukturę near-vertex następujący tekst po angielsku:

``The subscript $1$ records the distance parameter only. Both data structures accept every set $S\subseteq V$; their query bounds are stated as functions of $|S|$. In the application to \cref{lem:summary-dom}, the near-vertex oracle is queried only on sets of size at most $\ell$, whereas the far-vertex oracle is queried only on sets of size at most $k$. Thus these distance-$1$ structures provide a strictly stronger interface than the one required by that lemma.''}
\paragraph{Procedure \texorpdfstring{$\init{G}$ for $\nearVertexDS{G}{k,1}$}{Init for NearVertex[1,n]}.}
\begin{enumerate}
\item Initialize one instance of the deterministic toolkit on the edgeless graph on $V$.
\item For every edge $uv\in E(G)$, execute the insertion procedure of Lemma~\ref{lem:infra-update}.
\end{enumerate}
Because every intermediate graph is a subgraph of $G$, it is still $d$-degenerate, and because $G$ has at most $d(n-1)$ edges, this takes $\Oh(nd\cdot 2^{4d}(d^3+\log^3 n))$ time.

\paragraph{Procedure \texorpdfstring{$\add{uv}$ for $\nearVertexDS{G}{k,1}$}{Add for NearVertex[1,n]}.}
\begin{enumerate}
\item Execute the insertion procedure of Lemma~\ref{lem:infra-update} in the unique toolkit instance.
\end{enumerate}
This takes amortized time $\Oh(2^{4d}(d^3+\log^3 n))$.

\paragraph{Procedure \texorpdfstring{$\remove{uv}$ for $\nearVertexDS{G}{k,1}$}{Remove for NearVertex[1,n]}.}
\begin{enumerate}
\item Execute the deletion procedure of Lemma~\ref{lem:infra-update} in the unique toolkit instance.
\end{enumerate}
This takes amortized time $\Oh(2^{4d}(d^3+\log^3 n))$.

\paragraph{Procedure \texorpdfstring{$\near{S}$ for $\nearVertexDS{G}{k,1}$}{nearVertex(S) for NearVertex[1,n]}.}
\begin{enumerate}
\item If $S=\emptyset$, return vertex $1$.
\item Scan the vertices of $S$ and test whether they belong to $V^+$.
\item If some $s\in S$ is isolated, then return $s$ when $S=\{s\}$, and return $\bot$ otherwise.
\item For every $v\in S$, test whether $S\setminus\{v\}\subseteq N(v)$. If this holds for some $v$, return $v$.
\item Build the temporary de-duplicated list of candidates from $\bigcup_{s\in S}N^+(s)$.
\item For every candidate $u$ from Step~5, test whether $S\subseteq N(u)$. If this holds for some $u$, return $u$.
\item If $|S|\le \Delta$, inspect the bucket $B(S)$. If it is nonempty, return its first element.
\item Return $\bot$.
\end{enumerate}
Steps~2--3 take $\Oh(|S|)$ time, Step~4 takes $\Oh(|S|^2\log n)$ time, Steps~5--6 take $\Oh(d|S|\log n+d|S|^2\log n)$ time, and Step~7 takes $\Oh(d)$ time.
\end{proof}


\begin{theorem}\label{thm:distone-near}
Assume $n\ge 2$ and $d\ge 1$.
There exists a data structure $\nearVertexDS{G}{k,1}$ that maintains a dynamic $d$-degenerate graph $G$ on the fixed vertex set $V=\{1,\ldots,n\}$ and supports the following operations:
\begin{itemize}
\item $\init{G}$ initializes the structure in time $\Oh(nd\cdot 2^{4d}(d^3+\log^3 n))$;
\item $\add{uv}$ and $\remove{uv}$ update the structure in amortized time $\Oh(2^{4d}(d^3+\log^3 n))$;
\item $\near{S}$ returns a vertex $v$ with $S\subseteq N[v]$ if such a vertex exists, and returns $\bot$ otherwise, in time $\Oh(|S| + d|S|^2\log n)$.
\end{itemize}
The structure uses $\Oh(nd\cdot 2^{4d})$ words of space.
\end{theorem}
	
\end{comment}

\begin{proof}
The data structure consists of one instance of the toolkit of \cref{lem:infra-update} and auxiliary Boolean array indexed by $V$ used for temporary marking. Therefore, the guarantees on the update time, initialization time, and space usage follows directly from \cref{lem:infra-update}. We are left with implementing the query $\near{S,1}$.

% so the paragraphs above already define the full persistent state and all supported procedures.
%The bounds for $\init{G}$ and for the updates $\add{uv}$ and $\remove{uv}$ are exactly the ones stated after those procedures, and they follow directly from Lemma~\ref{lem:infra-update}.


Upon query $\near{S,1}$, we execute the following steps:
\begin{enumerate}
%	\item If $S=\emptyset$, return any vertex, say $1$.
	\item Initialize $L$ to be an empty list. The intention is that on $L$ we shall gather the vertices $u\in S\cup \bigcup_{s\in S}N^+(S)$ satisfying $S\subseteq N[u]$.
	\item For every $s\in S$, test whether $S\setminus\{s\}\subseteq N(s)$. If this is the case, add $s$ to $L$ and mark it in the auxiliary array as added.
	\item For every $s\in S$ and every vertex $u\in N^+(s)$, test whether $S\subseteq N(u)$. If this is the case and $u$ is not marked as already added, add $u$ to $L$ and mark it as added.
	\item Iterate through $L$ and remove all the markings from the auxiliary Boolean table. Thus it becomes empty (all-false) and $L$ contains no duplicates.
	\item Initialize $\gamma\coloneqq 0$. Iterate through $L$ and for every vertex $u\in L$ that does not satisfy $S\subseteq N^+(u)$, increment $\gamma$ by $\wei(u)$. Thus, we obtain $\gamma=\sum_{u\in L, S\nsubseteq N^+(u)} \wei(u)$.
	\item If $|S|\le \Delta$, find the bucket $B(S)$ in the dictionary of buckets, if existent.
	\item If $L$ is nonempty or $B(S)$ exists, return the first vertex of $L$ or the first vertex of $B(S)$; otherwise return $\bot$. For $\sum_{v\colon S\subseteq N[v]} \wei(v)$, return $\wei(B(S))+\gamma$, where $\wei(B(S))$ is replaced by $0$ if $B(S)$ is not present in the dictionary of buckets.
\end{enumerate}
Steps takes constant time. Step~2 takes $\Oh(|S|^2\log n)$ time, because it boils down to $\Oh(|S|^2)$ lookups in the neighborhood dictionaries. Step~3 takes $\Oh(d|S|^2\log n)$ time, because there are at most $d|S|$ vertices $u$ to consider, and for each of them we test $|S|$ vertices of $S$ for membership in $N(u)$, each in $\Oh(\log n)$ time. Step~4 takes $\Oh(|L|)\leq \Oh(d|S|)$ time. Step~5 takes $\Oh(d|S|^2\log n)$ time analogously to Step~3. Step~6 takes $\Oh(d+\log n)$ time, because this is the complexity of retrieving the bucket $B(S)$ from the dictionary of buckets. Finally, Step~7 takes constant time, because retrieving the weight and the first element of a bucket takes constant time. Thus, executing the query takes $\Oh(d|S|^2\log n)$ time in the worst case.

The correctness of the algorithm presented above follows immediately from \cref{cl:distone-near-candidates} and the observation that the contribution to the sum $\sum_{v\colon S\subseteq N[v]} \wei(v)$ of the vertices $v$ that do not satisfy $S\subseteq N^+(v)$ is exactly counted in the variable $\gamma$.
\end{proof}

\begin{comment}
It remains to analyze one query $\near{S}$.
Step~1 handles the case $S=\emptyset$ in $\Oh(1)$ time, so assume from now on that $S\neq\emptyset$.
Steps~2--3 inspect all vertices of $S$ and test in $\Oh(1)$ time per vertex whether they belong to $V^+$, hence together they take $\Oh(|S|)$ time.
If some $s\in S$ is isolated, then any vertex $v$ with $S\subseteq N[v]$ must equal $s$.
Consequently, Step~3 is correct: it returns $s$ when $S=\{s\}$ and returns $\bot$ when $|S|\ge 2$.

Assume now that every vertex of $S$ is non-isolated.
Any witness $v$ with $S\subseteq N[v]$ is either an internal candidate, meaning $v\in S$, or an external candidate, meaning $v\notin S$.
Step~4 tests every internal candidate explicitly.
To verify one such candidate it suffices to check whether every vertex of $S\setminus\{v\}$ belongs to $N(v)$.
This uses $|S|-1$ lookups in the neighbor dictionary of $v$, hence $\Oh(|S|\log n)$ time per candidate and $\Oh(|S|^2\log n)$ time in total for Step~4.

For external candidates we use Claim~\ref{cl:distone-near-candidates}.
It tells us that every external witness lies either in the explicit union $\bigcup_{s\in S}N^+(s)$ or, when $|S|\le \Delta$, in the single bucket $B(S)$.
Step~5 enumerates the union of out-neighborhoods by scanning the lists $N^+(s)$ for all $s\in S$; their total size is at most $\Delta|S|$.
We de-duplicate the resulting list using a temporary dictionary, so Step~5 costs $\Oh(\Delta|S|\log n)=\Oh(d|S|\log n)$.
Step~6 then verifies every explicit external candidate by checking the condition $S\subseteq N(u)$ through $|S|$ dictionary lookups, hence one candidate costs $\Oh(|S|\log n)$ and all explicit candidates together cost $\Oh(\Delta|S|^2\log n)=\Oh(d|S|^2\log n)$.

If $|S|\le \Delta$, then Step~7 inspects the bucket $B(S)$.
The lookup of that bucket by the explicit key $S$ costs $\Oh(|S|)\le \Oh(\Delta)=\Oh(d)$.
If the bucket is nonempty, then its first element is a correct witness and can be returned in $\Oh(1)$ time, because every vertex in $B(S)$ satisfies $S\subseteq N^+(u)\subseteq N(u)$.
If $|S|>\Delta$, Claim~\ref{lem:distone-U2-bucket} says that there is no bucket witness and Step~7 can be skipped.
Finally, if none of Steps~3, 4, 6, and 7 returns a vertex, then Claim~\ref{cl:distone-near-candidates} shows that no witness exists, so Step~8 correctly returns $\bot$.

The total query time is therefore $\Oh(|S| + |S|^2\log n + d|S|\log n + d|S|^2\log n + d)$, which simplifies to $\Oh(|S| + d|S|^2\log n)$, since $d\ge 1$, $|S|\ge 1$, and $\log n\ge 1$.

For the space bound, Lemma~\ref{lem:infra-update} gives $\Oh(nd\cdot 2^{4d})$ words for the unique toolkit instance, and the structure stores no additional persistent fields.
\end{comment}

\subsection{Data structure for \texorpdfstring{$\far{S,1}$}{far(S,1)}}\label{sec:distone-far}

In this section we give a data structure for the $\far{S,1}$ queries. Similarly to the reasoning presented in \cref{sec:far-be}, the idea is to first give a data structure that maintains weighted sums over vertices non-adjacent to $S$ (this part is entirely deterministic), and then lift it to an example-reporting data structure using \cref{lem:restore-by-sampling}. Therefore, our first goal is to prove the following.

\newcommand{\WeiFar}{\mathsf{WeightedFarSum}}
\newcommand{\sumFar}{\mathtt{sumFar}}

\begin{lemma}\label{lem:wei-sums-far}
	Fix $d\in \N$.
	Let $G$ be a dynamic graph on vertex set $V=\{1,\ldots,n\}$ that is $d$-degenerate at all times, and let $\wei\colon V\to \Z$ be a weight function fixed upon initialization. Then there exists a data structure $\WeiFar[G]$ for $G$ that supports the following query:
	\begin{itemize}
		\item $\sumFar(S)$: given $S\subseteq V$, return $\sum_{v\in V\setminus N[S]} \wei(v)$.
	\end{itemize}
	The amortized update time is $\Oh(2^{4d}\cdot d^2\cdot \log^2 n)$, and every query is answered in worst-case time $\Oh(2^{|S|}\cdot d|S|^2\cdot \log n)$. The data structure can be initialized for an edgeless $G$ in time $\Oh(n)$ and uses $\Oh(2^{4d}\cdot dn)$ space at all times.
\end{lemma}
\begin{proof}
	Observe that by the Inclusion-Exclusion Principle,
	\[\sum_{v\in V\setminus N[S]} \wei(v) = \sum_{v\colon N[v]\cap S=\emptyset} \wei(v)= \sum_{S'\subseteq S} (-1)^{|S'|}\cdot \sum_{v\colon S'\subseteq N[v]} \wei(v).\]
	Therefore, if we maintain one instance of the $\nearVertexDS{G}{}$ data structure provided by \cref{thm:NearVertex-new-deg}, then the value $\sum_{v\in V\setminus N[S]} \wei(v)$ can be computed using the formula above from the answers to $2^{|S|}$ queries to $\nearVertexDS{G}{}$: one for each subset $S'\subseteq S$. The complexity guarantees follow directly from the guarantees provided by \cref{thm:NearVertex-new-deg}.
\end{proof}

\michalin{Tekst poniżej jest oryginalny.}

The report procedure below will use inclusion-exclusion over closed neighborhoods.
We first record how to compute the weighted contribution of an intersection of closed neighborhoods using the same open-neighborhood decomposition as in the near-vertex structure.
For a nonempty set $T\subseteq V$ let
\[
U(T):=\{u\in V\setminus T : T\subseteq N(u)\},
\]
and split
\[
U_2(T):=\{u\in U(T): T\subseteq N^+(u)\},
\qquad
U_1(T):=U(T)\setminus U_2(T).
\]
Recall that $U_1(T)\subseteq \bigcup_{t\in T}N^+(t)$ and that $U_2(T)=B(T)$ when $|T|\le \Delta$, while $U_2(T)=\emptyset$ when $|T|>\Delta$.
The only additional issue for closed neighborhoods is that vertices of $T$ themselves may belong to the intersection.
This is handled by the following claim.

\begin{claim}\label{cl:distone-closed-decomp}
Let $T\subseteq V$ be nonempty and define
\[
C(T):=V^+\cap\bigcap_{t\in T}N_G[t],
\qquad
A(T):=\{a\in T\cap V^+ : T\setminus\{a\}\subseteq N(a)\}.
\]
Then the sets $A(T)$, $U_1(T)$, and $U_2(T)$ are pairwise disjoint and
\[
C(T)=A(T)\cup U_1(T)\cup U_2(T).
\]
\end{claim}
\begin{claimproof}[Proof.]
Let $u\in C(T)$.
If $u\in T$, then $u\in V^+$ and, for every $t\in T\setminus\{u\}$, the condition $u\in N_G[t]$ implies $t\in N(u)$.
Thus $u\in A(T)$.
If $u\notin T$, then $u\in V\setminus T$ and $T\subseteq N(u)$, so $u\in U(T)$ and hence $u\in U_1(T)\cup U_2(T)$.

Conversely, every vertex of $A(T)$ belongs to $V^+$ and to every closed neighborhood $N_G[t]$ with $t\in T$.
Every vertex of $U_1(T)\cup U_2(T)$ belongs to $U(T)$, hence is adjacent to every vertex of $T$ and has positive degree.
Thus all three sets are contained in $C(T)$.
Finally, $A(T)\subseteq T$, while $U_1(T)$ and $U_2(T)$ are contained in $V\setminus T$, and $U_1(T)\cap U_2(T)=\emptyset$ by definition.
Hence the three sets are pairwise disjoint.
\end{claimproof}

\paragraph{The predicate for \texorpdfstring{$\far{S}$}{far(S)}.}
Fix one auxiliary marker vertex $x\notin V$.
We assign the marker vertex the fixed identifier $\id(x):=n+1$; thus the fixed universe of $\widehat G$ is identified with $\{1,\ldots,n+1\}$ when Lemma~\ref{lem:restore-by-sampling} is applied.
Let $\widehat G$ be obtained from the current graph $G$ by adding $x$ and by possibly joining $x$ with an arbitrary finite set $P\subseteq V$.
Let
\[
P(\widehat G):=N_{\widehat G}(x)\cap V
\]
be the marker-neighborhood of $x$.
Outside the execution of a query the marker vertex is isolated, so $P(\widehat G)=\emptyset$.
During a query $\far{S}$ we temporarily set $P(\widehat G)=S$ by adding the edges $xs$ for all $s\in S$ and removing them after the retrieval query.
We use the convention $N_G[\emptyset]=\emptyset$.

Define a binary function $f$ on pairs $(\widehat G,u)$ by setting $f(\widehat G,x):=0$ and, for $u\in V$, setting
\[
f(\widehat G,u)=1
\quad\Longleftrightarrow\quad
u\in V^+\setminus N_G[P(\widehat G)].
\]
Thus, whenever $P(\widehat G)=S$, the positive vertices for $f$ are exactly the non-isolated vertices of $V\setminus N_G[S]$.
Isolated vertices are intentionally excluded from $f$ and will be handled separately in the construction of $\farVertexDS{G}{k,1}$.

\paragraph{The data structure \texorpdfstring{$\mathsf{WeightedSum}_{f,w}(\widehat G)$}{WeightedSum}.}
For a fixed weight function $w\colon V\cup\{x\}\to\Z$, assume that every value $w(u)$ and every maintained counter fits in $\Oh(1)$ machine words.
In the application through Lemma~\ref{lem:restore-by-sampling}, the only weights used are the sampled indicators and sampled identifier weights, so this word-size assumption is satisfied.
The structure consists of:
\begin{itemize}
\item one private instance of the deterministic toolkit of Section~\ref{sec:distone-infra} for the induced graph $G=\widehat G[V]$;
\item the current marker set $P=N_{\widehat G}(x)\cap V$, represented by a dynamic array, a Boolean membership array, and positions in the array, so insertions and deletions take $\Oh(1)$ time;
\item the global counter $W_w^{\mathrm{all}}:=\sum_{u\in V^+} w(u)$;
\item for every materialized bucket $B(Z)$, its header stores the counter $W_{w,Z}:=\sum_{u\in B(Z)} w(u)$.
\end{itemize}

\paragraph{The structure \texorpdfstring{$\mathsf{Verify}_{f}(\widehat G)$}{Verify}.}
We are in the binary case of Lemma~\ref{lem:restore-by-sampling}, because $f(\widehat G,u)\in\{0,1\}$ for every valid pair $(\widehat G,u)$.
Therefore no separate instance of $\mathsf{Verify}_{f}(\widehat G)$ is maintained.

\paragraph{Procedure \texorpdfstring{$\mathsf{WeightedSum}_{f,w}.\mathsf{Init}$}{WeightedSum.Init}.}
\begin{enumerate}
\item Initialize one private instance of the deterministic toolkit on the edgeless graph on $V$.
\item Set $P:=\emptyset$, set $W_w^{\mathrm{all}}:=0$, and note that no bucket header stores a weighted counter because no bucket is materialized yet.
\item For every edge $uv\in E(G)$, execute the edge-update procedure of $\mathsf{WeightedSum}_{f,w}$ for $\add{uv}$.
\item For every marker-neighbor $v\in N_{\widehat G}(x)\cap V$, execute the marker-update procedure for $\add{xv}$.
\end{enumerate}

\paragraph{Procedure \texorpdfstring{$\add{uv}$ and $\remove{uv}$ inside $\mathsf{WeightedSum}_{f,w}$}{WeightedSum edge update}.}
\begin{enumerate}
\item \begin{sloppypar}Execute the corresponding update in the private toolkit and read the resulting log $(L_{\mathrm{in}}^+,L_{\mathrm{out}}^+,\Lambda^+,\Lambda^-)$.\end{sloppypar}
\item For every vertex $u\in L_{\mathrm{out}}^+$, subtract $w(u)$ from $W_w^{\mathrm{all}}$.
\item For every vertex $u\in L_{\mathrm{in}}^+$, add $w(u)$ to $W_w^{\mathrm{all}}$.
\item For every deleted registration $(u,Z)\in \Lambda^-$, subtract $w(u)$ from the weighted counter field stored in the corresponding materialized bucket $B(Z)$; if this deletion empties the bucket, then that field disappears together with the bucket object.
\item For every added registration $(u,Z)\in \Lambda^+$, if the corresponding bucket $B(Z)$ has just become materialized, initialize its weighted counter field to $0$, and then add $w(u)$ to that field.
\end{enumerate}

\paragraph{Procedure \texorpdfstring{$\add{xv}$ and $\remove{xv}$ inside $\mathsf{WeightedSum}_{f,w}$}{WeightedSum marker update}.}
\begin{enumerate}
\item If the operation is $\add{xv}$, insert $v$ into $P$.
\item If the operation is $\remove{xv}$, remove $v$ from $P$.
\end{enumerate}
The private deterministic toolkit is not updated for edges incident with $x$, because it is maintained only for the induced graph $G=\widehat G[V]$.

\paragraph{Procedure \texorpdfstring{$\mathsf{WeightedSum}_{f,w}.\mathsf{Report}$}{WeightedSum.Report}.}
\begin{enumerate}
\item Let $P:=N_{\widehat G}(x)\cap V$.
\item If $P=\emptyset$, return $W_w^{\mathrm{all}}$.
\item Set $\mathsf{Ans}:=W_w^{\mathrm{all}}$.
\item For every nonempty subset $T\subseteq P$:
 \begin{enumerate}
 \item Set $E_T:=0$ and initialize an empty temporary set $\mathsf{seen}$.
 Enumerate the candidates from $T\cup\bigcup_{t\in T}N^+(t)$.
 For every enumerated candidate $u$ not yet in $\mathsf{seen}$, insert $u$ into $\mathsf{seen}$ and add $w(u)$ to $E_T$ if $u\in V^+$ and $T\subseteq N_G[u]$.
 Here $T\subseteq N_G[u]$ means that, for every $t\in T$, either $t=u$ or $t\in N(u)$, and these membership tests use the maintained neighbor dictionaries.
 \item Let $B_T$ be the stored weighted bucket counter $W_{w,T}$ if $|T|\le \Delta$ and $B(T)$ is materialized; otherwise let $B_T:=0$.
 \item Set $C_T:=E_T+B_T$.
 \item If $|T|$ is odd, set $\mathsf{Ans}:=\mathsf{Ans}-C_T$; if $|T|$ is even, set $\mathsf{Ans}:=\mathsf{Ans}+C_T$.
 \end{enumerate}
\item Return $\mathsf{Ans}$.
\end{enumerate}

\begin{lemma}[Weighted sums for the hard part of \texorpdfstring{$\far{S}$}{far(S)}]\label{lem:distone-weighted-sum}
Assume $n\ge 2$ and $d\ge 1$.
For any fixed weight function $w\colon V\cup\{x\}\to\Z$, there exists a data structure $\mathsf{WeightedSum}_{f,w}(\widehat G)$ that maintains a dynamic graph $\widehat G$ of the form described above and, after each update, reports $\sum_{u\in V\cup\{x\}} f(\widehat G,u)w(u)$.
\begin{itemize}
\item initialization takes time $\Oh(nd\cdot 2^{4d}(d^3+\log^3 n))$;
\item every edge insertion or deletion inside the induced graph $G=\widehat G[V]$ is handled in amortized time $\Oh(2^{4d}(d^3+\log^3 n))$;
\item insertion or deletion of an edge incident with the marker $x$ is handled in amortized time $\Oh(1)$;
\item $\mathsf{WeightedSum}_{f,w}.\mathsf{Report}$ takes amortized time $\Oh(1)$ when $P=\emptyset$, and $\Oh(2^s\cdot d\cdot s^2\log n)$ when $s:=|P|\ge 1$.
\end{itemize}
The structure uses $\Oh(nd\cdot 2^{4d})$ words of space besides the read-only representation of $w$.
\end{lemma}
\begin{proof}
The private toolkit contributes $\Oh(1)$ initialization time on the edgeless graph, $\Oh(2^{4d}(d^3+\log^3 n))$ amortized time per edge insertion or deletion inside $G$, and $\Oh(nd\cdot 2^{4d})$ words of space by Lemma~\ref{lem:infra-update}.
It remains to justify the additional fields and the report formula.

In $\mathsf{WeightedSum}_{f,w}.\mathsf{Init}$, Step~1 initializes the private toolkit.
Step~2 initializes the marker set and the global counter, and, because the graph is still edgeless, leaves all bucket headers absent.
Step~3 then inserts the edges of $G$ one by one by calling the edge-update procedure of $\mathsf{WeightedSum}_{f,w}$, while Step~4 handles all current edges incident with $x$.
Because $G$ has at most $d(n-1)=\Oh(nd)$ edges and every marker update costs $\Oh(1)$ time, the total initialization time is $\Oh(nd\cdot 2^{4d}(d^3+\log^3 n))$.

Consider one edge insertion or deletion inside $G$.
Step~1 of the edge-update procedure updates the private toolkit and outputs the log $(L_{\mathrm{in}}^+,L_{\mathrm{out}}^+,\Lambda^+,\Lambda^-)$.
By Lemma~\ref{lem:infra-update}, the total number of logged items is $\Oh(2^{4d}(d+\log n))$ in amortized time.
Steps~2--5 perform one $\Oh(1)$ arithmetic operation per logged item.
Indeed, whenever a registration event is read from the toolkit log, the corresponding bucket object has already been identified in Step~1, so updating its weighted counter field costs $\Oh(1)$; if a bucket becomes materialized, its new field is initialized once, and if a bucket becomes empty, that field disappears together with the bucket object.
Hence the counter maintenance takes $\Oh(2^\Delta(d+\log n))$ amortized time, which is absorbed by the amortized cost $\Oh(2^{4d}(d^3+\log^3 n))$ of Step~1.

If the updated edge is incident with the marker $x$, then the edge belongs to $\widehat G$ but not to the induced graph $G=\widehat G[V]$.
In this case the private toolkit does not change, and the marker-update procedure only modifies the current marker set $P$.
The array representation described above supports this in $\Oh(1)$ time.

It remains to justify $\mathsf{WeightedSum}_{f,w}.\mathsf{Report}$.
If $P=\emptyset$, then $N_G[P]=\emptyset$, so the positive vertices of $f$ are exactly $V^+$ and Step~2 correctly returns $W_w^{\mathrm{all}}$.
Assume therefore that $P\neq\emptyset$.
By inclusion-exclusion,
\[
\sum_{u\in V\cup\{x\}} f(\widehat G,u)w(u)
=
W_w^{\mathrm{all}}
+
\sum_{\emptyset\neq T\subseteq P}(-1)^{|T|}
W(C(T)),
\]
where $W(C(T)):=\sum_{u\in C(T)}w(u)$ and $C(T)=V^+\cap\bigcap_{t\in T}N_G[t]$.

Fix one nonempty $T\subseteq P$.
By Claim~\ref{cl:distone-closed-decomp}, the sets $A(T)$, $U_1(T)$, and $U_2(T)$ are pairwise disjoint and
\[
C(T)=A(T)\cup U_1(T)\cup U_2(T).
\]
The explicit enumeration in Step~4a contains $A(T)$ because $A(T)\subseteq T$, and it contains $U_1(T)$ by Claim~\ref{lem:distone-U1-bound}.
The test $u\in V^+$ and $T\subseteq N_G[u]$ then accepts exactly the enumerated vertices of $C(T)$.
The remaining part is $U_2(T)$.
By Claim~\ref{lem:distone-U2-bucket}, this is the bucket $B(T)$ when $|T|\le \Delta$, and is empty when $|T|>\Delta$.
Thus $B_T$ is exactly the weighted contribution of $U_2(T)$.
Moreover, the bucket contribution is disjoint from the explicit part: if $u\in B(T)$, then $T\subseteq N^+(u)$, so for every $t\in T$ the edge $ut$ is oriented from $u$ to $t$; hence $u\notin N^+(t)$, and also $u\notin T$.
Therefore $u$ is never enumerated from $T\cup\bigcup_{t\in T}N^+(t)$.
Consequently $C_T=E_T+B_T=W(C(T))$, and the signs in Step~4d implement the displayed inclusion-exclusion formula.

For the running time, the case $P=\emptyset$ takes $\Oh(1)$ time.
Let $s:=|P|\ge 1$.
There are $2^s-1$ nonempty subsets $T\subseteq P$.
For one such $T$, Step~4a enumerates at most $|T|+\Delta|T|=\Oh(d|T|)$ candidates.
For each candidate it checks $T\subseteq N_G[u]$ using $\Oh(|T|)$ neighbor-dictionary lookups, hence $\Oh(|T|\log n)$ time per candidate.
The bucket lookup costs $\Oh(|T|)\le \Oh(d|T|)$ when $|T|\le \Delta$ and is skipped otherwise.
Thus the work for one $T$ is $\Oh(d|T|^2\log n)$, and the whole report time is $\Oh(2^s\cdot d\cdot s^2\log n)$.

For the space bound, the private toolkit uses $\Oh(nd\cdot 2^{4d})$ words.
The counter $W_w^{\mathrm{all}}$ contributes $\Oh(1)$ words, and the marker set representation contributes $\Oh(n)$ words.
Finally, each materialized bucket header stores one additional weighted counter field.
The number of materialized buckets is at most the total number of registrations plus one, and every vertex has at most $2^\Delta$ registrations.
Hence all additional weighted counter fields together use $\Oh(n\cdot 2^\Delta)=\Oh(nd\cdot 2^{4d})$ words.
Since $d\ge 1$, the marker-set storage is absorbed, and the whole structure uses $\Oh(nd\cdot 2^{4d})$ words.
\end{proof}

\paragraph{The data structure \texorpdfstring{$\farVertexDS{G}{k,1}$}{FarVertex[1,n]}.}
The structure consists of:
\begin{itemize}
\item one auxiliary marker vertex $x\notin V$, kept isolated outside the execution of a query;
\item one private instance of the deterministic toolkit of Section~\ref{sec:distone-infra}, used only to inspect isolated vertices of $G$;
\item one randomized data structure $\mathsf{RetrieveVertex}_{f}(\widehat G)$ obtained from Lemma~\ref{lem:restore-by-sampling} applied black-box to the full structure $\mathsf{WeightedSum}_{f,w}(\widehat G)$.
\end{itemize}

\paragraph{Procedure \texorpdfstring{$\init{G}$ for $\farVertexDS{G}{k,1}$}{Init for FarVertex[1,n]}.}
\begin{enumerate}
\item Initialize one private instance of the deterministic toolkit on the edgeless graph on $V$.
\item For every edge $uv\in E(G)$, execute the insertion procedure of Lemma~\ref{lem:infra-update} in this toolkit instance.
\item Keep the marker vertex $x$ isolated and form the graph $\widehat G:=G+x$.
\item Initialize $\mathsf{RetrieveVertex}_{f}(\widehat G)$ on $\widehat G$.
\end{enumerate}

\bbinline{Analiza inicjalizacji struktury $\farVertexDS{G}{1}$ odtwarza krawędzie dowolnego grafu początkowego, podczas gdy \cref{thm:main-domset-deg} wymaga inicjalizacji na grafie pustym. Proponuję dopisać bezpośrednio po wyświetlonej procedurze inicjalizacji $\farVertexDS{G}{1}$ następujący akapit po angielsku:

``When the initial graph is edgeless, Step~1 initializes the private toolkit in $\Oh(1)$ time and Step~2 performs no work. The weighted-sum structures used in Step~4 are likewise initialized on an edgeless graph. Consequently, \cref{lem:restore-by-sampling} yields initialization time $\Oh(n\log n\log(1/\eps))$ for one instance of $\mathsf{RetrieveVertex}_{f}(\widehat G)$. Hence one instance of $\farVertexDS{G}{1}$ can be initialized on the edgeless graph in $\Oh(n\log n\log(1/\eps))$ time. In the proof of \cref{thm:main-domset-deg}, this bound is multiplied by the number of independent far-vertex instances maintained there.''}

\paragraph{Procedure \texorpdfstring{$\add{uv}$ for $\farVertexDS{G}{k,1}$}{Add for FarVertex[1,n]}.}
\begin{enumerate}
\item Relay the insertion independently to the private toolkit and to $\mathsf{RetrieveVertex}_{f}(\widehat G)$.
\end{enumerate}

\paragraph{Procedure \texorpdfstring{$\remove{uv}$ for $\farVertexDS{G}{k,1}$}{Remove for FarVertex[1,n]}.}
\begin{enumerate}
\item Relay the deletion independently to the private toolkit and to $\mathsf{RetrieveVertex}_{f}(\widehat G)$.
\end{enumerate}

\paragraph{Procedure \texorpdfstring{$\far{S}$ for $\farVertexDS{G}{k,1}$}{farVertex(S) for FarVertex[1,n]}.}
\begin{enumerate}
\item Canonicalize $S$ as a set and build a temporary membership dictionary for it.
\item Query the interval dictionary of the private toolkit for the smallest isolated vertex $z$.
\item While such a vertex $z$ exists and $z\in S$, query the interval dictionary for the next isolated vertex after $z$.
\item If the resulting isolated vertex $z$ exists, return it immediately.
\item Temporarily execute $\add{xs}$ on $\mathsf{RetrieveVertex}_{f}(\widehat G)$ for every $s\in S$.
\item Query $\mathsf{RetrieveVertex}_{f}(\widehat G)$.
\item Temporarily execute $\remove{xs}$ on $\mathsf{RetrieveVertex}_{f}(\widehat G)$ for every $s\in S$.
\item If the retrieval structure reported a vertex, return it; otherwise return $\bot$.
\end{enumerate}

\paragraph{Randomness model.}
The random choices inside Lemma~\ref{lem:restore-by-sampling} are fixed during the initialization of $\mathsf{RetrieveVertex}_{f}(\widehat G)$.
The adversary is oblivious to these choices with respect to the entire operation sequence seen by this randomized retrieval structure, including the temporary marker updates $\add{xs}$ and $\remove{xs}$ for all $s\in S$ performed inside a query $\far{S}$.

\bbinline{Bieżący argument losowy nie wystarcza w zastosowaniu przez semi-ladder algorithm: kolejne zapytania far-vertex są adaptacyjnie wyznaczane na podstawie wcześniejszych odpowiedzi. Proponuję zastąpić aktualny akapit rozpoczynający się od dosłownego nagłówka "Randomness model" pierwszym z poniższych tekstów po angielsku. Drugi tekst warto zachować jako regułę stosowania tej struktury w dowodzie \cref{thm:main-domset-deg}; precyzuje on niezależne instancje dla rund i argument dla całego czasu życia każdej instancji.

``The random choices inside \cref{lem:restore-by-sampling} are fixed when $\mathsf{RetrieveVertex}_{f}(\widehat G)$ is initialized. Against an oblivious adversary, the guarantee applies to every fixed operation sequence seen by this retrieval structure, including every temporary marker-edge insertion $\add{xs}$ and deletion $\remove{xs}$ performed inside a call to $\far{S}$.''

``When this primitive is used by the semi-ladder algorithm, do not reuse one randomized far-vertex instance in all rounds. Instead, maintain independent instances $\mathcal F_1,\ldots,\mathcal F_\ell$, where $\mathcal F_i$ is used only in round $i$. Fix the complete external edge-update sequence. For each $i$, condition on the random bits of $\mathcal F_1,\ldots,\mathcal F_{i-1}$. The complete operation sequence issued to $\mathcal F_i$ throughout the lifetime of the data structure is then determined by the external updates, the deterministic near-vertex oracle, and the earlier instances, and is independent of the random bits of $\mathcal F_i$. Therefore the oblivious-adversary guarantee applies to every round-$i$ far-vertex query. If every $\mathcal F_i$ has error parameter $\delta$, then a union bound over at most $\ell$ rounds gives error probability at most $\ell\delta$ for one semi-ladder execution.''}

\begin{theorem}\label{thm:distone-far}
Assume $n\ge 2$, $d\ge 1$, and let $\eps \in (0,1/2)$.
There exists a randomized data structure $\farVertexDS{G}{k,1}$ that, against an oblivious adversary, maintains a dynamic $d$-degenerate graph $G$ on the fixed vertex set $V=\{1,\ldots,n\}$ and supports the following operations:
\begin{itemize}
\item $\init{G}$ initializes the structure in time
$\Oh\!\left(nd\cdot 2^{4d}(d^3+\log^3 n)\log n \log \frac{1}{\eps}\right)$;
\item $\add{uv}$ and $\remove{uv}$ update the structure in amortized time
$\Oh\!\left(2^{4d}(d^3+\log^3 n)\log n \log \frac{1}{\eps}\right)$;
\item $\far{S}$, for any $S\subseteq V$, behaves as follows:
 \begin{itemize}
 \item if $V\setminus N_G[S]=\emptyset$, it returns $\bot$;
 \item if $V\setminus N_G[S]\neq\emptyset$, it never returns a vertex of $N_G[S]$, and it returns some vertex of $V\setminus N_G[S]$ with probability at least $1-\eps$.
 \end{itemize}
\end{itemize}
The amortized query time on input $S$, where $s:=\max(1,|S|)$, is
\[
\Oh\!\left((s+2^s d s^2\log n)\log n \log \frac{1}{\eps}\right).
\]
The structure uses $\Oh\!\left(nd\cdot 2^{4d}\log n \log \frac{1}{\eps}\right)$ words of space.
\end{theorem}

\begin{proof}
The predicate $f$ was defined above so that, after temporarily adding the marker edges $xs$ for all $s\in S$, the condition $f(\widehat G,u)=1$ holds exactly for the non-isolated vertices $u\in V\setminus N_G[S]$.
Thus the only witnesses not represented by $f$ are isolated vertices, and these are handled by the private toolkit stored directly in $\farVertexDS{G}{k,1}$.
An isolated vertex $z$ is a valid witness for $\far{S}$ exactly when $z\notin S$.

By Lemma~\ref{lem:distone-weighted-sum}, one full instance of $\mathsf{WeightedSum}_{f,w}(\widehat G)$ has initialization time
$I_R=\Oh(nd\cdot 2^{4d}(d^3+\log^3 n))$,
amortized update time
$T_R=\Oh(2^{4d}(d^3+\log^3 n))$,
amortized report time
$Q_R(s)=\Oh(2^s\cdot d\cdot s^2\log n)$ for $s:=\max(1,|P|)$,
and space
$S_R=\Oh(nd\cdot 2^{4d})$.
Because $f$ is binary, we apply Lemma~\ref{lem:restore-by-sampling} in the case \cref{it:binary-f}, hence $I_V=T_V=Q_V=S_V=0$.
Since $|\widehat G|=n+1=\Oh(n)$, $d\ge 1$, and $\eps\in(0,1/2)$, this yields a randomized data structure $\mathsf{RetrieveVertex}_{f}(\widehat G)$ with initialization time $\Oh(nd\cdot 2^{4d}(d^3+\log^3 n)\log n \log \frac{1}{\eps})$, amortized update time $\Oh(2^{4d}(d^3+\log^3 n)\log n \log \frac{1}{\eps})$, amortized query time $\Oh(Q_R(s)\log n \log \frac{1}{\eps})$ when the current marker set has size at most $s$, and space $\Oh(nd\cdot 2^{4d}\log n \log \frac{1}{\eps})$.

We now construct $\farVertexDS{G}{k,1}$ from one private toolkit instance together with $\mathsf{RetrieveVertex}_{f}(\widehat G)$.
By the randomness paragraph above, the oblivious-adversary guarantee applies to the entire operation sequence seen by $\mathsf{RetrieveVertex}_{f}(\widehat G)$, including the temporary marker updates inside a call to $\far{S}$.
In $\init{G}$, Steps~1--2 initialize the private toolkit on the input graph in time $\Oh(nd\cdot 2^{4d}(d^3+\log^3 n))$ by Lemma~\ref{lem:infra-update}.
Step~3 is constant time and starts with $P(\widehat G)=\emptyset$, and Step~4 initializes $\mathsf{RetrieveVertex}_{f}(\widehat G)$ with the bound stated above.
This gives the claimed initialization time.

For one later update $\add{uv}$ or $\remove{uv}$, the graph edge $uv$ is relayed to the private toolkit and to $\mathsf{RetrieveVertex}_{f}(\widehat G)$.
The private toolkit contributes amortized time $\Oh(2^{4d}(d^3+\log^3 n))$ by Lemma~\ref{lem:infra-update}, while the black-box structure $\mathsf{RetrieveVertex}_{f}(\widehat G)$ contributes amortized time $\Oh(2^{4d}(d^3+\log^3 n)\log n \log \frac{1}{\eps})$.
Marker updates do not occur in ordinary $\add{uv}$ and $\remove{uv}$ operations; they are performed only temporarily inside queries.
This proves the claimed update bound.

We next analyze one call to $\far{S}$.
The canonicalization and temporary membership dictionary for $S$ take $\Oh(|S|\log n)$ time.
Steps~2--4 inspect isolated vertices using the interval dictionary of the private toolkit.
The loop can skip only isolated vertices that belong to $S$, so it iterates at most $|S|$ times and costs $\Oh((|S|+1)\log n)$ time.
If this stage returns an isolated vertex $z$, then $z\notin S$, and since $z$ has no neighbors, $z\notin N_G[S]$.
Thus the returned vertex is a correct witness.

Assume now that Steps~2--4 do not return.
Then no isolated vertex outside $S$ exists, so every remaining witness in $V\setminus N_G[S]$ is non-isolated.
After Step~5, the current marker set is $P(\widehat G)=S$, and by the definition of $f$ the positive vertices for $f$ are exactly $V^+\setminus N_G[S]$.
Since every remaining witness is non-isolated, these positive vertices are exactly the remaining witnesses for $\far{S}$.
Step~5 performs $|S|$ marker-edge insertions and Step~7 performs $|S|$ marker-edge deletions.
Each such update costs $\Oh(\log n \log \frac{1}{\eps})$ through $\mathsf{RetrieveVertex}_{f}(\widehat G)$, because the corresponding update in $\mathsf{WeightedSum}_{f,w}$ costs $\Oh(1)$.
Step~6 is one retrieval query and, while $P(\widehat G)=S$, costs
\[
\Oh\!\left(2^{|S|}d|S|^2\log^2 n \log \frac{1}{\eps}\right)
\]
for $|S|\ge 1$; the case $S=\emptyset$ is covered by using $s=\max(1,|S|)$.
Together with the isolated-vertex stage, this gives the displayed query bound.

If Step~6 reports a vertex, then by the no-false-positives guarantee of Lemma~\ref{lem:restore-by-sampling} this vertex satisfies $f(\widehat G,u)>0$, hence $u\in V\setminus N_G[S]$.
If Step~6 reports no vertex, then Step~8 returns $\bot$.
When $V\setminus N_G[S]=\emptyset$, the isolated-vertex stage cannot find a valid isolated witness and the retrieval stage has no positive vertex to return.
When $V\setminus N_G[S]\neq\emptyset$, the isolated-vertex stage finds a witness whenever an isolated witness exists; otherwise the witness set is exactly the set of positive vertices of $f$ after Step~5, so Lemma~\ref{lem:restore-by-sampling} guarantees that Step~6 returns some witness with probability at least $1-\eps$.
This proves the one-sided error model stated in the lemma.

Finally, the space of $\farVertexDS{G}{k,1}$ is the sum of the space of one private toolkit instance, namely $\Oh(nd\cdot 2^{4d})$, and the space of $\mathsf{RetrieveVertex}_{f}(\widehat G)$, namely $\Oh(nd\cdot 2^{4d}\log n \log \frac{1}{\eps})$.
The marker set $P$ inside the weighted structures contributes only $\Oh(n)$ words per instance and is absorbed in the same bound because $d\ge 1$.
Hence the total space is $\Oh(nd\cdot 2^{4d}\log n \log \frac{1}{\eps})$.
\end{proof}

\bbinline{Obecne granice w \cref{thm:main-domset-deg} nie wynikają z \cref{lem:summary-dom}: koszt generowania kandydatów daje narzut $2^{k^{\Oh(d)}}$, a niezależne instancje far-vertex wprowadzają dodatkowy czynnik zależny od liczby rund. Proponuję zastąpić trzy końcowe zdania opisujące złożoność w \cref{thm:main-domset-deg} pierwszym z poniższych tekstów po angielsku. Następnie proponuję wstawić bezpośrednio w tym miejscu drugi tekst — pełny dowód twierdzenia. Tekst ten jawnie obsługuje $k=0$, $n\le1$, $d=0$ i $k=1$, dobiera poprawnie parametr błędu oraz korzysta tylko z deterministycznego rdzenia \cref{lem:summary-dom}.

``Put $\bar\eps\coloneqq\min\{\eps,1/4\}$. Every answer to a query is correct with probability at least $1-\eps$ against an oblivious adversary. The amortized time complexity of an update or a query is
\[
 2^{\Oh(d)+k^{\Oh(d)}}\log^4 n\left(\log\frac{1}{\bar\eps}+d\log(k+1)\right).
\]
The data structure can be initialized on an edgeless graph in time
\[
 (k+1)^{\Oh(d)}2^{\Oh(d)}n\log n\left(\log\frac{1}{\bar\eps}+d\log(k+1)\right)
\]
and occupies the same asymptotic amount of space at all times.''

}

\bbinline{``We first dispose of the degenerate parameter cases. If $k=0$, return $\emptyset$ if and only if $V(G)=\emptyset$, and return $\bot$ otherwise. If $n\le 1$ or $d=0$, then $G$ is edgeless at all times; in this case return $V(G)$ if $|V(G)|\le k$, and return $\bot$ otherwise. Hence assume that $n\ge 2$, $d\ge 1$, and $k\ge 1$.

Let $\mathcal C_d$ be the class of all $d$-degenerate graphs. Every graph in $\mathcal C_d$ excludes $K_{d+1,d+1}$ as a subgraph, and therefore \cref{thm:sli-biclique} gives $\sli_1(\mathcal C_d)<3(d+1)$. Put
\[
 \ell\coloneqq\max\{3(d+1),k^{3(d+1)}\},\qquad
 \bar\eps\coloneqq\min\{\eps,1/4\},\qquad
 \delta\coloneqq\bar\eps/\ell.
\]
For $k\ge2$, \cref{thm:termination} implies that every semi-ladder execution has fewer than $k^{3(d+1)}$ rounds. For $k=1$, suppose that an execution has $q$ rounds. In its first $q-1$ rounds the algorithm receives witnesses $w_1,\ldots,w_{q-1}$, and the corresponding candidates are singletons $D_i=\{a_i\}$. The pairs $(a_i,w_i)$ form a distance-$1$ semi-ladder of order $q-1$. Hence $q\le\sli_1(\mathcal C_d)+1\le3(d+1)$. Thus every execution has at most $\ell$ rounds.

}

\bbinline{Use the deterministic near-vertex oracle from \cref{thm:distone-near}; although it accepts arbitrary sets, it is queried only on sets of size at most $\ell$. For the far-vertex oracle, maintain $\ell$ independent instances $\mathcal F_1,\ldots,\mathcal F_\ell$ of \cref{thm:distone-far}, each initialized with error parameter $\delta$, and relay every edge update to all of them. In round $i$ of every semi-ladder execution, answer $\notDominated{D_i,1}$ by querying $\mathcal F_i.\far{D_i}$; here $|D_i|\le k$.

Fix the complete external edge-update sequence. For a fixed $i$, condition on all random bits of $\mathcal F_1,\ldots,\mathcal F_{i-1}$. Throughout the lifetime of the data structure, the complete operation sequence issued to $\mathcal F_i$, including all temporary marker-edge insertions and deletions, is determined by the external updates, the deterministic near-vertex oracle, and the earlier instances. In particular, it is independent of the random bits of $\mathcal F_i$. Therefore the oblivious-adversary guarantee of \cref{thm:distone-far} applies to every round-$i$ query, which fails with probability at most $\delta$. A union bound over the at most $\ell$ rounds of the current semi-ladder execution gives failure probability at most $\ell\delta=\bar\eps\le\eps$.

Conditional on all far-oracle calls made in the current semi-ladder execution being correct, the deterministic argument in the proof of \cref{lem:summary-dom}, together with \cref{lem:candidate-to-near}, proves correctness of the returned answer. This uses only the deterministic core of \cref{lem:summary-dom}; its randomized `Moreover' paragraph is not invoked. The near-vertex oracle contributes no error.

}

\bbinline{The $\ell$ far-vertex instances multiply their update time, initialization time, and space by $\ell$, and
\[
 \log\frac1\delta=\log\frac1{\bar\eps}+\log\ell
 =\log\frac1{\bar\eps}+\Oh(d\log(k+1)).
\]
Together with the $2^{k^{\Oh(d)}}$ candidate-query overhead from \cref{lem:summary-dom}, this yields update time
\[
 2^{\Oh(d)+k^{\Oh(d)}}\log^4 n\left(\log\frac1{\bar\eps}+d\log(k+1)\right).
\]
The edgeless-graph initialization bound stated above for one far-vertex instance gives total initialization time and space
\[
 (k+1)^{\Oh(d)}2^{\Oh(d)}n\log n\left(\log\frac1{\bar\eps}+d\log(k+1)\right).
\]
This proves the theorem.''}
